\documentclass[10pt,a4paper]{article}

\usepackage[margin=20mm]{geometry}
\usepackage{amsmath,amssymb,amsthm}

\theoremstyle{definition}
\newtheorem*{problem*}{Problem}

\setlength{\parindent}{0pt}
\setlength{\parskip}{5pt}

\begin{document}

\begin{center}
  {\Large\bfseries The Sharp Curvature-Oscillation Radius}
\end{center}

\section{Definitions}

For a smooth regular immersed closed planar curve $\gamma$, let $L$ be
its length, $A$ its signed area, $k$ its signed curvature, and $s$ its
arclength parameter. Define
\begin{equation}
  N(\gamma)
  =
  \frac{1}{2\pi}\int_\gamma k\,ds,
  \qquad
  \overline k
  =
  \frac{1}{L}\int_\gamma k\,ds,
\end{equation}
and
\begin{equation}
  K_{\mathrm{osc}}(\gamma)
  =
  L\int_\gamma(k-\overline k)^2\,ds.
\end{equation}
For a curve $\gamma_0$ satisfying $N(\gamma_0)=1$, let
\begin{equation}
  T_{\max}(\gamma_0)\in(0,\infty]
\end{equation}
be its maximal smooth existence time under curve diffusion, with the sign
chosen so that length decreases. Define
\begin{equation}
  \varepsilon_{\mathrm{CD}}
  =
  \sup\left\{
    \varepsilon>0:
    K_{\mathrm{osc}}(\gamma_0)<\varepsilon
    \Longrightarrow
    T_{\max}(\gamma_0)=\infty
    \text{ for every such }\gamma_0
  \right\}.
\end{equation}

\section{Best Known Result}

The sharp inequality
\begin{equation}
  1-\frac{4\pi A}{L^2}
  \leq
  \frac{K_{\mathrm{osc}}}{8\pi^2}
\end{equation}
shows that sufficiently small curvature oscillation forces $A>0$.
Existing small-data theory therefore gives
$\varepsilon_{\mathrm{CD}}>0$.
Conversely, any smooth rotation-one curve with $A\leq 0$ has finite maximal
time, so such curves give a finite upper bound for
$\varepsilon_{\mathrm{CD}}$.

\section{Problem Formalization}

Determine the exact value of $\varepsilon_{\mathrm{CD}}$.

An exact value is ideal. Until it is known, a valid result must provide a
theorem-certified bracket
\begin{equation}
  L\leq\varepsilon_{\mathrm{CD}}\leq U.
\end{equation}
Unsupported point estimates are not sufficient.

\end{document}